% Cote 119: offprints, reports, courses, CV and letters (typed and printed part): archive cote 119, site page 294 (PDF page 295).
% typescript. Draft compared: none.
% AI-transcribed by Claude Opus 5.5 (claude-opus-5-5, adaptive thinking, effort high), Message Batches API, batch r119_full (msgbatch_01DwZNm5Y7YsKBWyAN4sFs8w), 2026-11-13; not yet checked by eye.
\scanpage{119/294}
\mspage{- 7 -}

définis; je ne sais même pas si les applications
\[ i_{\hand{\rho}} : H^1(k, E_8) \rightarrow H^3(k, Z/2Z) \]
du n°3.2 peuvent être non triviales.

\textbf{§4. Problèmes d'injectivité}

L'ensemble $H^1(k, G)$ est fonctoriel en $k$ et $G$ :

a) Si $k'$ est une extension de $k$, on a une application naturelle
\[ H^1(k, G) \rightarrow H^1(k', G) . \]
b) Si $G \rightarrow G'$ est un morphisme de groupes algébriques, on a une
application naturelle $H^1(k, G) \rightarrow H^1(k, G')$.

On dispose d'une série de cas où ces applications sont \uline{injectives} :

(4.1\unsure{\hand{)}}- \uline{théorème de simplification de Witt}) Si $q = q_1 \hand{\oplus} q_2$ , où les $q_i$
sont des formes quadratiques, l'application $H^1(k, O(q_1)) \rightarrow H^1(k, O(q))$
est injective.

(4.2)- Même énoncé, pour les \uline{groupes unitaires} associés aux algèbres à involution
sur $k$.

Ce résultat, nettement plus délicat que le précédent, a fait l'objet d'un
exposé par E.Bayer.

(4.3) (Springer)- Injectivité de $H^1(k, O(q)) \rightarrow H^1(k', O(q))$ lorsque
$k'$ est une extension finie de $k$ \uline{de degré impair}.

(4.4) (Bayer-Lenstra)-Même énoncé que (4.3), pour les \uline{groupes unitaires}
au lieu des groupes orthogonaux.

(4.5) (Pfister)-Injectivité de $H^1(k, O(q)) \rightarrow H^1(k, O(q \hand{\otimes} q'))$ lorsque
le rang de $q'$ est impair (le morphisme $O(q) \rightarrow O(q \hand{\otimes} q')$ étant défini
par le produit tensoriel).

On aimerait avoir d'autres énoncés du même type, par exemple les suivants
(qui sont peut-être trop optimistes) :

(4.6 ?)-Si $k'$ est une extension finie de $k$ de degré premier à 2 et
3, l'application $H^1(k, F_4) \rightarrow H^1(k', F_4)$ est injective.

(4.7 ?)- Même énoncé pour $E_8$ , avec $\hand{\{}2,3\hand{\}}$ remplacé par $\hand{\{}2,3,5\hand{\}}$ .

\uline{Remarque} - Soit $G$ un groupe algébrique sur $k$, et soient $x,y$ deux éléments

\margin{294}
